\documentclass[12pt,a4paper]{article}

\usepackage{styles/couleurs}
\usepackage{styles/mise_en_page}
\usepackage{styles/boites}
\usepackage{styles/maths}
\usepackage{styles/figures}

\usepackage{enumitem}
\usepackage{array}
\usepackage{tabularx}
\usepackage{amssymb}
\usepackage{tikz}

\newcommand{\lignes}[1]{%
  \begin{tcolorbox}[colback=white,colframe=NeutreC!40,boxrule=0.5pt,arc=1.5mm,left=3mm,right=3mm,top=2mm,bottom=1mm]
    \foreach \i in {1,...,#1}{\noindent\dotfill\par\vspace{0.45em}}
  \end{tcolorbox}
}

\begin{document}

\begin{center}
  {\sffamily\LARGE\bfseries\color{TitreC}Séquence 1 -- Statistiques}\\[0.2em]
  {\sffamily\large\bfseries\color{SavoirC}Activité 3 -- Défi : médiane 5}
\end{center}

\remarque{Dans cette activité, il ne s'agit pas de calculer une médiane à partir d'une série donnée : c'est à toi de construire des séries qui respectent une contrainte.}

\section*{C'est quoi déjà?}
\textbf{À quoi sait-on qu'une série a bien pour médiane \(5\) ?}
\lignes{4}

\section*{Défi A}
Construis une série de \textbf{sept nombres tous différents} dont la médiane est \(5\).
\lignes{3}

\section*{Défi B}
Construis une série de \textbf{six nombres} dont la médiane est \(5\) et dont \textbf{au moins un nombre est égal à \(5\)}.
\lignes{3}

\section*{Défi C}
Construis une série de \textbf{huit nombres, tous différents de \(5\)}, dont la médiane est \(5\).
\lignes{4}

% NOTES PROFESSEUR -- NON IMPRIMÉES
% Plusieurs réponses possibles. Validation possible par échange entre binômes.
% Relances : « Comment peux-tu vérifier ta réponse ? » ;
% « Que se passe-t-il une fois la série rangée ? » ;
% « Dans le défi C, la médiane doit-elle forcément appartenir à la série ? »

\end{document}
